% \documentclass[letterpaper]{article}
% \usepackage[submission]{aaai2027}
% \usepackage[hyphens]{url}
% \usepackage{graphicx}
% \urlstyle{rm}
% \def\UrlFont{\rm}
% \usepackage{natbib}
% \usepackage{caption}
% \frenchspacing
% \usepackage{booktabs}
% \usepackage{multirow}
% \usepackage{amsmath}
% \usepackage{amssymb}
% \usepackage{xcolor}
% \usepackage{array}

% \newcommand{\ours}{ZipTok3D}
% \pdfinfo{/TemplateVersion (2027.1)}

% \renewcommand{\thesection}{S\arabic{section}}
% \renewcommand{\thesubsection}{S\arabic{section}.\arabic{subsection}}
% \renewcommand{\thefigure}{\Alph{figure}}
% \renewcommand{\thetable}{\Alph{table}}
% \renewcommand{\theequation}{S\arabic{equation}}
% \setcounter{secnumdepth}{2}

% \title{Supplementary Material for \ours{}: High-Fidelity 3D Tokenization with Compact Token Prefixes}
% \author{Anonymous Submission}
% \affiliations{}

% \begin{document}
% \maketitle

% \section{Prefix-Length and Refinement-Depth Analysis}
% \label{sec:full-sweeps}

% \subsection{Complete Metric Sweeps}

% We evaluate prefix lengths
% $K\in\{1,2,4,8,16,32,64,128\}$ and refinement depths
% $L\in\{1,2,3,4,5,6\}$ using one selected checkpoint per dataset. At every
% operating point, the decoder receives a $K\times512$ latent tensor without
% suffix padding.

% \begin{figure}[!t]
%     \centering
%     \begin{minipage}{\columnwidth}
%         \centering
%         \includegraphics[width=.79\columnwidth]
%         {supplementary_figures/mesh_f1_interactions.pdf}
%         \captionof{figure}{Effect of prefix length $K$ and refinement depth
%         $L$ on mesh F1@0.02 for ShapeNet (top) and TRELLIS (bottom). The plots
%         vary $K$ at fixed $L$ (left) and $L$ at fixed $K$ (right).}
%         \label{fig:supp-f1-curves}
%     \end{minipage}
%     \begin{minipage}{\columnwidth}
%         \centering
%         \includegraphics[width=.79\columnwidth]
%         {supplementary_figures/mesh_cd_interactions.pdf}
%         \captionof{figure}{Effect of prefix length $K$ and refinement depth
%         $L$ on mesh CD for ShapeNet (top) and TRELLIS (bottom). The plots vary
%         $K$ at fixed $L$ (left) and $L$ at fixed $K$ (right). Lower values are
%         better.}
%         \label{fig:supp-cd-curves}
%     \end{minipage}
% \end{figure}

% Figures~\ref{fig:supp-f1-curves} and~\ref{fig:supp-cd-curves} extend the
% query-IoU analysis to the two surface metrics. Both exhibit the same asymmetric
% interaction: increasing $K$ is most beneficial under shallow decoding, whereas
% increasing $L$ is most beneficial for short prefixes. The
% differences between five and six passes are small at the reported operating
% points across query IoU, mesh CD, and mesh F1, supporting the use of five
% passes for evaluation.

% Figure~\ref{fig:supp-heatmaps} reports the complete grids for $L=2,\ldots,6$.
% The single-pass surface-metric cross-sections, whose larger ranges would
% compress the heatmap colors, are reported separately in
% Figures~\ref{fig:supp-f1-curves} and~\ref{fig:supp-cd-curves}.

% \begin{figure*}[!t]
%     \centering
%     {\setlength{\tabcolsep}{1.5pt}
%     \begin{tabular}{@{}cc@{}}
%     \includegraphics[width=.48\textwidth]{supplementary_figures/shapenet_query_iou.pdf} &
%     \includegraphics[width=.48\textwidth]{supplementary_figures/trellis_query_iou.pdf} \\[-.55em]
%     \includegraphics[width=.48\textwidth]{supplementary_figures/shapenet_mesh_f1_0p02.pdf} &
%     \includegraphics[width=.48\textwidth]{supplementary_figures/trellis_mesh_f1_0p02.pdf} \\[-.55em]
%     \includegraphics[width=.48\textwidth]{supplementary_figures/shapenet_mesh_cd.pdf} &
%     \includegraphics[width=.48\textwidth]{supplementary_figures/trellis_mesh_cd.pdf}
%     \end{tabular}}
%     \caption{Complete ShapeNet (left) and TRELLIS (right) sweeps for
%     refinement depths $L=2,\ldots,6$ and all supported prefix lengths. Rows
%     report query IoU, mesh F1@0.02, and mesh CD. Cell values are dataset means.
%     Darker cells indicate better performance; the color direction is reversed
%     for CD.}
%     \label{fig:supp-heatmaps}
% \end{figure*}

% \subsection{Per-Object Refinement Behavior}

% \begin{table*}[!t]
% \centering
% \small
% \setlength{\tabcolsep}{5pt}
% \begin{tabular}{lccccc}
% \toprule
% Dataset & $K$ & Transition & IoU improved (\%) & CD improved (\%) & F1 improved (\%) \\
% \midrule
% \multirow{4}{*}{ShapeNet}
%  & 1 & $1\!\rightarrow\!3$ & 100.0 & 99.3 & 97.6 \\
%  & 1 & $3\!\rightarrow\!5$ & 88.2 & 81.1 & 58.1 \\
%  & 4 & $1\!\rightarrow\!3$ & 100.0 & 98.7 & 94.8 \\
%  & 4 & $3\!\rightarrow\!5$ & 89.5 & 81.5 & 55.7 \\
% \midrule
% \multirow{4}{*}{TRELLIS}
%  & 1 & $1\!\rightarrow\!3$ & 98.6 & 99.8 & 99.8 \\
%  & 1 & $3\!\rightarrow\!5$ & 91.2 & 91.8 & 84.7 \\
%  & 4 & $1\!\rightarrow\!3$ & 98.5 & 99.7 & 99.4 \\
%  & 4 & $3\!\rightarrow\!5$ & 90.5 & 90.7 & 83.4 \\
% \bottomrule
% \end{tabular}
% \caption{Per-object refinement diagnostics. ``Improved'' means a strict
% increase for IoU/F1 or a strict decrease for CD between the indicated passes.
% The remaining objects are ties or regressions.}
% \label{tab:supp-refinement-diagnostics}
% \end{table*}

% Mean improvements do not imply monotonic improvement for every object.
% Table~\ref{tab:supp-refinement-diagnostics} therefore reports strict
% per-object improvement rates between selected refinement depths. Query
% comparisons include all 1,283 ShapeNet and 2,613 TRELLIS objects. Mesh
% comparisons exclude seven
% TRELLIS objects for which at least one reconstructed mesh in the compared pair
% is empty or invalid, leaving 2,606 paired objects; all ShapeNet meshes are
% valid.

% Increasing depth from $L=1$ to $L=3$ improves nearly every valid object, so the
% aggregate gain is not driven by a small subset. Increasing depth from $L=3$ to
% $L=5$ remains beneficial for a majority of objects but is not uniformly
% monotonic, particularly for ShapeNet mesh F1. Later passes therefore yield
% smaller aggregate gains and may slightly degrade individual reconstructions.

% \section{Experimental Setup Details}
% \label{sec:experimental-details}

% \subsection{Datasets, Preprocessing, and Splits}

% \paragraph{ShapeNet.}
% We use ShapeNetCore-v2 \citep{chang2015shapenet} with the 55-category
% preprocessing and split manifests released by 3DShape2VecSet
% \citep{zhang20233dshape2vecset} and adopted by COD-VAE
% \citep{cho2025representing}.  The manifests contain 48,597 training, 1,283
% validation, and 2,592 test objects.  For each object, the released data provide
% uniform volume queries, near-surface queries, their occupancy labels, a
% normalization scale, and a dense surface point cloud.  Following COD-VAE, we
% consolidate the records by split while preserving category and object
% identifiers.  All manifest entries used in the reported experiments are
% successfully converted.
% Query coordinates are stored in FP16 for training and FP32 otherwise, while
% surface points remain in FP32.  Training samples 2,048 surface points as
% encoder input and 4,096 queries from each of the volume and near-surface sets.
% During training, independent axis scales are sampled uniformly from
% $[0.75,1.25]$, the shape is renormalized to the unit cube, and Gaussian noise
% with standard deviation 0.005 is added to the surface points before clipping
% to $[-1,1]$.  The same geometric scaling is applied to the occupancy queries.
% Validation and test data are not augmented, and evaluation uses all stored
% volume queries and the full surface point cloud. The reported reconstruction,
% ablation, and refinement analyses use the 1,283-object split named
% \emph{validation} in the released manifests. The 2,592-object test split is
% retained as part of the released partition but is not used in the reported
% tokenizer tables.

% \paragraph{TRELLIS.}
% TRELLIS-500K \citep{xiang2025structured} provides metadata for 500,777
% training assets: 168,307 Objaverse-XL Sketchfab assets, 311,843 Objaverse-XL
% GitHub assets, 4,485 ABO assets, 9,472 3D-FUTURE assets, and 6,670 HSSD assets.
% Its separate 3,229-object Toys4K evaluation set is not used.  Our preprocessing
% starts from the available downloaded assets rather than the complete metadata
% catalogue.  We merge supported polygonal components and apply the
% point-cloud-utils watertight conversion used by the VecSet preprocessing
% pipeline at resolution 50,000 with seed 0.  The watertight mesh is centered at
% its bounding-box center and scaled so that its maximum vertex radius is one.
% Each successful object is represented by 100,000 surface points, 500,000
% volume queries sampled uniformly from the radius-$\sqrt{3}$ sphere, and
% 500,000 near-surface queries divided equally between Gaussian perturbation
% scales 0.005 and 0.05. Occupancy is one where the signed distance is
% nonpositive and zero otherwise.
% Per-object sampling is deterministic from seed 42 and the object identifier.
% Inputs above 50,000 faces are decimated to that target before watertight
% conversion when decimation succeeds. The preprocessing additionally sets a
% 900-second timeout and a 32-GiB memory limit per worker. An asset is excluded
% if it lacks a supported,
% nonempty polygonal mesh, contains nonfinite geometry, or fails mesh conversion,
% normalization, surface sampling, signed-distance evaluation, or the stated
% resource limits.

% The accompanying manifest fixes the held-out evaluation split at 2,613
% objects: 57 ABO,
% 856 Objaverse-XL GitHub, and 1,700 Objaverse-XL Sketchfab assets.  The original
% metadata catalogue is therefore not the effective split pool. With seed 42,
% 1\% of the successfully preprocessed pool is assigned to validation, and the
% remaining non-evaluation objects are assigned to training. The manifest
% contains the source category and exact object identifier of each evaluation
% asset and defines the TRELLIS split used for all reconstruction, bootstrap,
% and adaptive-budget analyses.

% \subsection{Baseline Sources and Operating Points}

% We compare with 3DILG \citep{zhang20223dilg}, VecSet
% \citep{zhang20233dshape2vecset}, and COD-VAE
% \citep{cho2025representing}.  These methods support deterministic
% reconstruction at fixed, shape-independent token budgets.  All ShapeNet
% comparisons use the common 1,283-object split and COD-VAE evaluation pipeline.
% The ShapeNet baseline implementations and checkpoints are inherited from the
% official COD-VAE release. On TRELLIS, COD-VAE-2/32/64 are initialized from
% their corresponding ShapeNet checkpoints and fine-tuned on the common
% preprocessed training pool. VecSet-512 uses the inherited ShapeNet model
% without TRELLIS-specific training. TRELLIS results are not reported for
% 3DILG or the shorter VecSet operating points. All reported TRELLIS methods
% use the common split and evaluation protocol described above.

% The reconstruction comparisons use 512 tokens for 3DILG; 32, 64, and 512
% tokens for VecSet; and 2, 32, and 64 tokens for COD-VAE.  ZipTok3D is evaluated
% with exact prefixes of 1, 2, and 4 tokens.  The class-conditioned generation
% comparison uses the native stage-2 latent lengths of the baselines and the
% exact two-token stage-2 representation of ZipTok3D.

% \subsection{Optimization and Model Selection}
% \label{sec:optimization}

% \begin{table*}[t]
% \centering
% \small
% \setlength{\tabcolsep}{4pt}
% \begin{tabular}{lcc}
% \toprule
% Setting & ShapeNet & TRELLIS \\
% \midrule
% Optimizer & \multicolumn{2}{c}{AdamW, $\beta_1=0.9$, $\beta_2=0.999$,
% $\epsilon=10^{-8}$, weight decay 0.01} \\
% Initial learning rate & \multicolumn{2}{c}{$10^{-4}$} \\
% Precision / clipping & \multicolumn{2}{c}{FP16 mixed precision / global-norm
% threshold 0.5} \\
% Batching & \multicolumn{2}{c}{56 examples per GPU, four A800 GPUs, three-step
% gradient accumulation; effective batch size 672} \\
% Training duration / schedule & 1,000 epochs; 50-epoch warmup, then cosine decay
% & 300 epochs; constant learning rate \\
% Seed / repetitions & \multicolumn{2}{c}{123456 / one run per configuration} \\
% Checkpoint selection & \multicolumn{2}{c}{Highest mean query IoU on the
% corresponding validation split} \\
% \bottomrule
% \end{tabular}
% \caption{Stage-1 optimization and model-selection protocol.}
% \label{tab:supp-training}
% \end{table*}

% Each training example contains 2,048 surface points, 4,096 uniformly sampled
% volume queries, and 4,096 near-surface queries.  The reconstruction criterion
% weights the volume and near-surface terms by 1.0 and 0.1, respectively.  The
% intermediate-loss coefficient is $\beta=0.5$.  At optimizer step $t$,
% $\alpha_t=\max(0,1-t/T_\alpha)$, with $T_\alpha$ equal to the scheduled number
% of optimizer updates after distributed batching and gradient accumulation.
% Stage-1 training unrolls six refinement passes and samples one intermediate
% depth uniformly from $\{1,\ldots,5\}$.  Thus $\alpha_t$ decays over the full
% training run and reaches zero at its final scheduled update.
% The implementation uses Python 3.9, PyTorch 2.1.0, and CUDA 12.2.

% \subsection{Stage-2 Generation Training}

% The prefix VAE uses AdamW with learning rate $10^{-4}$,
% weight decay 0.01, global batch size 376, FP16 mixed precision, and gradient
% clipping at 0.5.  It is trained for 100 epochs, with the learning rate halved at
% epochs 60, 70, 80, and 90.  The EDM uses AdamW with learning rate $10^{-4}$,
% weight decay 0.05, global batch size 256, FP16 mixed precision, and gradient
% clipping at 1.0.  Its 1,000-epoch schedule uses a 40-epoch linear warmup followed
% by cosine decay to $10^{-6}$.  Both stages use four A800 GPUs and seed 123456.
% For EDM training we set $p_{\mathrm{mean}}=-1.2$,
% $p_{\mathrm{std}}=1.2$, and $\sigma_{\mathrm{data}}=1$.  We select the
% prefix VAE checkpoint by the highest validation query IoU and the EDM
% checkpoint by the lowest validation denoising loss.

% \section{Architecture and Loss Specifications}
% \label{sec:architecture}

% \subsection{Stage-1 Tokenizer Architecture}

% Table~\ref{tab:supp-architecture} specifies the stage-1 architecture. Unless
% noted otherwise, these components follow the released COD-VAE design
% \citep{cho2025representing}. ZipTok3D introduces nested prefix training, the
% shared refinement block, and intermediate supervision.

% \begin{table*}[t]
% \centering
% \small
% \setlength{\tabcolsep}{5pt}
% \begin{tabular}{p{.19\textwidth}p{.74\textwidth}}
% \toprule
% Component & Configuration \\
% \midrule
% Input and point encoder & 2,048 surface points; 512 farthest-point-sampled
% patches; width 512; four progressive encoder blocks with three self-attention
% layers per block; eight attention heads; GEGLU MLP ratio 4; dropout 0.1. \\
% Latent bank & Maximum length $M=128$ and width 512.  The initial latent queries
% are point embeddings of 128 farthest-point-sampled surface positions; no
% separate learned latent-position table is used. \\
% Prefix sampling & One budget is sampled uniformly from
% $\{1,2,4,8,16,32,64,128\}$ for each training example.  All suffix positions are
% masked from every decoder attention operation. \\
% Initial triplane tokens & Three planes, 32 channels, spatial resolution
% $128\times128$, and patch size $8\times8$, giving 768 learned spatial tokens
% $T^0$. \\
% Selection and merging & Within $\mathcal{S}_{\omega}$, one cross-attention
% layer lets $T^0$ interact with the retained prefix. A LayerNorm--MLP
% uncertainty head then ranks the 768 spatial tokens. The top 25\% (192 tokens)
% form $H_K^0$; the remaining tokens are summarized by cross-attention into eight
% merged tokens $R_K$. \\
% Shared refinement & One six-layer Transformer block of width 512, eight heads,
% GEGLU MLP ratio 2, and dropout 0.1 is reused for every pass. Each call receives
% the current 192 selected tokens and the retained $K$-token prefix. Only the
% selected-token output is carried to the next pass. The eight merged tokens
% $R_K$ bypass the shared block and remain fixed for restoration. There is no
% causal mask, pass embedding, or loop-specific parameter. \\
% Restoration and triplane head & Updated selected tokens are scattered back to
% their original spatial positions.  A linear residual head maps every spatial
% token to an $8\times8\times32$ patch and adds it to the initial triplane patch,
% modulated by the predicted uncertainty. \\
% Occupancy head & Features bilinearly sampled from the three planes are summed
% and passed through a $32\!\rightarrow\!32\!\rightarrow\!1$ MLP with GELU. \\
% \bottomrule
% \end{tabular}
% \caption{Detailed stage-1 tokenizer architecture.}
% \label{tab:supp-architecture}
% \end{table*}

% \subsection{Shared Refinement and Auxiliary Losses}

% The recurrent sequence at pass $\ell$ can therefore be written as
% \begin{equation}
% U_K^{\ell}=
% \bigl[H_K^{\ell-1};Z_{:K}\bigr]
% \in\mathbb{R}^{(192+K)\times512},
% \end{equation}
% where $H_K^{\ell-1}$ is the selected state and $Z_{:K}$ is the retained prefix.
% The shared block processes this sequence, but only its first 192 outputs define
% $H_K^{\ell}$. The prefix is reinserted unchanged at every pass, while the
% merged representation $R_K$ remains outside the block and is used only when
% restoring the complete triplane.

% The two inherited COD-VAE auxiliary losses stabilize triplane initialization
% and uncertainty-based selection. Let $a_j^0$ be the initial occupancy logit for
% query $j$, $\hat y_j^0=\sigma(a_j^0)$ its probability, and
% $\hat{\mathcal{Y}}^0=(\hat y_j^0)_{j=1}^{N_q}$ the corresponding probability
% vector. The initialization loss applies the reconstruction criterion to this
% initial prediction:
% \begin{equation}
% \mathcal{L}_{\mathrm{init}}
% =\mathcal{L}_{\mathrm{rec}}(\hat{\mathcal{Y}}^0,\mathcal{Y}).
% \end{equation}
% The uncertainty target is the detached pointwise binary cross-entropy of the
% initial prediction.  Following COD-VAE, it is clipped and normalized to
% $[0,1]$ before regression:
% \begin{align}
% r_j & =
% \frac{\operatorname{clip}(\operatorname{BCE}(\hat y_j^0,y_j)-0.01,0,0.99)}{0.99},\\
% \mathcal{L}_{\mathrm{unc}}
% &=\frac{1}{N_q}\sum_{j=1}^{N_q}(u_j-\operatorname{sg}(r_j))^2,
% \end{align}
% where $u_j$ is the uncertainty field decoded at the query. We use
% $\lambda_{\mathrm{init}}=0.5$ and $\lambda_{\mathrm{unc}}=0.001$.

% \subsection{Stage-2 Prefix VAE and EDM}

% The prefix VAE accepts up to 16 ordered stage-1 tokens.  A LayerNorm and linear
% projection independently map each 512-dimensional token to the mean and log
% variance of a 32-dimensional diagonal Gaussian posterior.  The log variance is
% clamped to $[-30,20]$, and training samples from the posterior by the
% reparameterization trick.  A linear projection maps the sampled sequence back
% to width 512, after which learned positional embeddings and a 12-layer
% Transformer reconstruct the stage-1 token features.  This Transformer uses
% eight heads, GEGLU MLP ratio 4, and dropout 0.1.  Its causal mask is
% left-to-right:
% \begin{equation}
% A_{ij}=\begin{cases}
% 0, & j\leq i,\\
% -\infty, & j>i,
% \end{cases}
% \end{equation}
% so the reconstruction at position $i$ cannot use a later token.

% Let $\widetilde Z$ denote a posterior sample and
% $\bar Z=D_{\psi}(\widetilde Z)$ the reconstructed stage-1 token sequence.  The
% stage-1 tokenizer is frozen while the prefix VAE is jointly optimized for
% $\mathcal{B}=\{1,2,4,8,16\}$ with weights
% $w_K\in\{1.2,1.1,1.0,0.9,0.8\}$.  Its objective is
% \begin{equation}
% \begin{aligned}
% \mathcal{L}_{\mathrm{pVAE}}
% &=\sum_{K\in\mathcal{B}}
% \frac{w_K}{\sum_{K'\in\mathcal{B}}w_{K'}}
% \\[-0.2em]
% &\quad\times\Bigl[
% \operatorname{MSE}\!\left(\bar Z_{:K},\operatorname{LN}(Z_{:K})\right)
% \\[-0.2em]
% &\hspace{3.2em}
% +\mathcal{L}_{\mathrm{rec}}\!\left(F_6(\bar Z_{:K}),\mathcal{Y}\right)
% \Bigr]
% \\
% &\quad+10^{-3}D_{\mathrm{KL}}\!\left(
% q_{\psi}(\widetilde Z\mid Z_{:16})\,\|\,\mathcal{N}(0,I)
% \right).
% \end{aligned}
% \end{equation}
% where $F_6$ denotes the frozen stage-1 decoder and occupancy head evaluated
% with six refinement passes.  The feature and occupancy terms both have unit
% coefficients; the volume and near-surface weighting inside
% $\mathcal{L}_{\mathrm{rec}}$ is given in Section~\ref{sec:optimization}.
% Causality ensures that every term for budget $K$ depends only on the
% corresponding stage-2 prefix. The six-pass decoder is used only to supervise
% prefix VAE training. During generation evaluation, sampled codes are decoded
% using five refinement passes.

% We cache deterministic posterior means for the full $16\times32$ latent array,
% normalize them using training-set channel statistics, and train the
% class-conditional EDM on these arrays. Left-to-right causal self-attention
% ensures that the denoising prediction at each position depends only on the
% corresponding noisy prefix. At the reported two-token operating point,
% sampling therefore instantiates only the first two latent positions, and the
% prefix VAE decoder receives an exact $2\times32$ tensor without suffix padding.
% The EDM follows the latent-array design used by
% VecSet and COD-VAE \citep{zhang20233dshape2vecset,cho2025representing}: it has
% width 384, 16 Transformer blocks, eight 48-dimensional heads, GEGLU MLP ratio
% 4, no dropout, learned token-position embeddings, and a left-to-right causal
% self-attention mask. A
% 256-dimensional sinusoidal noise embedding is projected to width 384 and
% provides scale and shift to the adaptive LayerNorms in every block.  The class
% index spans all 55 ShapeNet categories. It is mapped to a learned
% 384-dimensional embedding and supplied as a one-token key/value context
% through cross-attention in every block; token self-attention therefore models
% the latent sequence while class conditioning is injected separately.

% For a normalized latent array $x$, category $c$, and
% $\log\sigma\sim\mathcal{N}(p_{\mathrm{mean}},p_{\mathrm{std}}^2)$, the EDM
% objective is
% \begin{equation}
% \begin{aligned}
% \mathcal{L}_{\mathrm{EDM}}
% &=\mathbb{E}_{x,c,\sigma,\epsilon}\!\left[
% \frac{\sigma^2+\sigma_{\mathrm{data}}^2}
% {(\sigma\sigma_{\mathrm{data}})^2}
% \right.\\[-0.2em]
% &\hspace{4.2em}\left.
% \times\left\|D_{\theta}(x+\epsilon;\sigma,c)-x\right\|_2^2
% \right],
% \\
% &\hspace{4.2em}
% \epsilon\sim\mathcal{N}(0,\sigma^2I).
% \end{aligned}
% \end{equation}
% All latent positions present at an operating point are denoised in parallel;
% there is no token-autoregressive sampling loop.

% \section{Evaluation Protocols}
% \label{sec:evaluation}

% \subsection{Reconstruction Metric Settings}

% Query IoU uses all 500,000 stored volume queries per object on both datasets.
% For mesh evaluation, we query a dense
% $128^3$ occupancy grid, extract the zero-logit level set, equivalently the
% 0.5-probability isosurface, with
% marching cubes, and sample 100,000 points from both the reconstruction and
% reference surfaces. Chamfer distance is the sum of the two mean bidirectional
% Euclidean nearest-neighbor distances. At threshold 0.02, precision and recall are
% the fractions of reconstructed and reference points, respectively, lying
% within the threshold of the other surface.  F1 is their harmonic mean.

% \subsection{Class-Conditioned Generation Metric Settings}

% Following COD-VAE \citep{cho2025representing}, we evaluate airplane, car,
% chair, table, and rifle and generate 2,000 shapes per category.  We report
% minimum matching distance (MMD-CD), coverage (COV-CD), and 1-nearest-neighbor
% accuracy (1-NNA-CD). For each category, the reference set $S_r$ is the
% corresponding subset of the final ShapeNet evaluation split. From the 2,000
% generated shapes, we use $5|S_r|$ for MMD-CD and COV-CD and $|S_r|$ for
% 1-NNA-CD. Each mesh is represented
% by 2,048 points sampled from its surface, and category-level scores are averaged
% over the five categories.  Using CD as the set distance, MMD averages the
% distance from each reference shape to its nearest generated shape.  COV is the
% fraction of reference shapes selected as a nearest neighbor of at least one
% generated shape.  1-NNA is the leave-one-out nearest-neighbor classification
% accuracy on the union of equally sized reference and generated sets.

% VecSet and COD-VAE use 18 diffusion steps, whereas 3DILG uses its native
% 512-step autoregressive sampler.  Our EDM uses 18 Heun steps with $\rho=7$,
% $\sigma_{\max}=80$, and $\sigma_{\min}=0.002$, without condition dropout or
% classifier-free guidance.

% \subsection{Efficiency Benchmark Protocol}

% All efficiency measurements use the corresponding trained checkpoints on one
% NVIDIA H20 GPU with batch size 16 and FP32 inference.  TF32 and Flash SDP are
% disabled.  CUDA events measure GPU execution after warmup; model loading,
% data loading, CPU--GPU transfers, marching cubes, and file output are excluded.
% For reconstruction, full inference begins with the input point cloud and
% includes point encoding, tokenizer decoding, and all $128^3$ occupancy queries
% in chunks of 131,072 points. Decoder latency instead begins with an already
% available stage-1 prefix and ends when the complete triplane has been
% produced; it excludes point encoding and occupancy queries.

% For generation, sampling throughput measures EDM latent generation alone.
% Full throughput begins from the same latent noise and additionally includes
% prefix VAE decoding, five-pass stage-1 tokenizer decoding, and the complete
% $128^3$ occupancy-field query. Thus the point encoder is part of reconstruction
% full inference but not of the generation pipeline. We discard three warmup
% batches and average ten measured batches for reconstruction and decoder
% latency; generation uses five warmup and thirty measured batches. For a
% measured mean batch latency $t$ in milliseconds,
% \begin{equation}
% \mathrm{throughput}=\frac{16\times1000}{t}.
% \end{equation}
% Peak allocated CUDA memory is measured after resetting the peak-memory counter
% following warmup. It is measured over the corresponding full pipeline;
% reserved memory is not used.

% \begin{table}[t]
% \centering
% \small
% \setlength{\tabcolsep}{4pt}
% \begin{tabular}{lrrr}
% \toprule
% Method & Tok. & Full shape/s & Peak GiB \\
% \midrule
% VecSet & 512 & 4.20 & 32.63 \\
% COD-VAE & 2 & 82.48 & 2.08 \\
% COD-VAE & 32 & 79.82 & 2.08 \\
% COD-VAE & 64 & 77.29 & 2.09 \\
% \midrule
% \ours{} & 1 & 72.51 & 2.03 \\
% \ours{} & 2 & 72.51 & 2.03 \\
% \ours{} & 4 & 72.24 & 2.03 \\
% \bottomrule
% \end{tabular}
% \caption{Full reconstruction efficiency under the common protocol.  ZipTok3D
% uses five refinement passes.}
% \label{tab:supp-reconstruction-efficiency}
% \end{table}

% Table~\ref{tab:supp-reconstruction-efficiency} reports the resulting
% representation--compute tradeoff. Shorter prefixes reduce neither the dense
% occupancy query nor the number of shared refinement calls. Consequently,
% ZipTok3D has lower end-to-end throughput than the single-pass COD-VAE variants,
% despite its shorter latent sequence. Peak memory remains similar because the
% common encoder and dense-field query dominate the allocation.

% \subsection{Statistical Analysis}

% For the primary comparison between ZipTok3D at $K=4,L=5$ and COD-VAE-32, we
% resample aligned TRELLIS objects with replacement. Each of 20,000 resamples
% contains 2,613 object indices. For each metric, we compute the difference
% between the paired sample means and report the 2.5th and 97.5th percentiles of
% the resulting bootstrap distribution. The bootstrap uses seed 123456, and all
% effects are defined as ZipTok3D minus COD-VAE, so negative CD is favorable.
% These intervals quantify variation over evaluation objects conditional on the
% trained models; they do not measure variation across independent training
% runs.

% \section{Post-Hoc Adaptive-Budget Diagnostic}
% \label{sec:posthoc-oracle}

% To characterize how the required budget varies across individual shapes, we
% perform a post-hoc diagnostic over
% $K\in\{1,2,4,8,16,32\}$ and $L\in\{1,\ldots,6\}$ from the same checkpoint.
% Candidates are searched by increasing $K$ and then increasing $L$, prioritizing
% a shorter representation before shallower decoding.  This diagnostic accesses
% ground-truth reconstruction metrics and thus measures adaptive potential rather
% than the performance of a deployable budget predictor.

% We use a strict criterion that accepts a candidate only when
% all three metrics are no worse than COD-VAE-32 after decimal half-up rounding
% (one decimal for IoU/F1 and three decimals for CD).  Shapes without a successful
% candidate are excluded from this diagnostic.  Table~\ref{tab:supp-posthoc-all}
% reports both methods on the same successful subset, preventing differences in
% subset composition from affecting the comparison.

% \begin{table}[t]
% \centering
% {\small\setlength{\tabcolsep}{2pt}
% \begin{tabular*}{\columnwidth}{@{\extracolsep{\fill}}llcccc@{}}
% \toprule
% Data & Method & Found & Rate & Avg. $K$ & Avg. $L$ \\
% \midrule
% ShapeNet & COD-VAE-32 & -- & -- & 32.00 & -- \\
% ShapeNet & \ours{} & 979/1,283 & 76.31 & 2.43 & 3.41 \\
% TRELLIS & COD-VAE-32 & -- & -- & 32.00 & -- \\
% TRELLIS & \ours{} & 1,767/2,613 & 67.62 & 3.02 & 3.97 \\
% \bottomrule
% \end{tabular*}
% \vspace{0.35em}

% \begin{tabular*}{\columnwidth}{@{\extracolsep{\fill}}llccc@{}}
% \toprule
% Data & Method & IoU$\uparrow$ & CD$\downarrow$ & F1$\uparrow$ \\
% \midrule
% ShapeNet & COD-VAE-32 & 97.2 & 0.011 & 98.1 \\
% ShapeNet & \ours{} & 97.3 & 0.011 & 98.3 \\
% TRELLIS & COD-VAE-32 & 75.8 & 0.015 & 96.1 \\
% TRELLIS & \ours{} & 76.6 & 0.014 & 96.6 \\
% \bottomrule
% \end{tabular*}}
% \caption{All-metric post-hoc diagnostic. The \ours{} rows report the post-hoc
% oracle. Found and Rate give the number and percentage of shapes for which the
% search identifies an operating point that matches or exceeds COD-VAE-32 on all
% three rounded per-shape metrics. Avg. $K$ and Avg. $L$ are computed over these
% shapes. Metrics in the lower block are aggregated over the same subset for
% both methods, so the COD-VAE rows are subset rather than full-split results.
% IoU and F1 are percentages.}
% \label{tab:supp-posthoc-all}
% \end{table}

% No candidate in the evaluated grid satisfies the strict all-metric criterion
% for the remaining 304 ShapeNet and 846 TRELLIS objects. All fixed-budget
% evaluations include these objects; they are excluded only from the oracle
% averages in Table~\ref{tab:supp-posthoc-all}. Consequently, the
% reported average budgets characterize the successful subset rather than the
% expected cost of a deployable early-exit policy.

% Among successful objects, three or four passes account for 68.3\% and 71.6\%
% of the ShapeNet and TRELLIS subsets, respectively. This distribution is
% consistent with the dataset-level saturation behavior in
% Figures~\ref{fig:supp-f1-curves} and~\ref{fig:supp-cd-curves}.

% \begin{figure}[t]
% \centering
% \includegraphics[width=\columnwidth]{supplementary_figures/adaptive_success_token_distribution.pdf}
% \caption{First successful prefix length under the token-first all-metric
% search.  Percentages are normalized within the successful subset of each
% dataset.}
% \label{fig:supp-posthoc-token-dist}
% \end{figure}

% \begin{figure}[t]
% \centering
% \includegraphics[width=\columnwidth]{supplementary_figures/adaptive_success_loop_distribution.pdf}
% \caption{First successful refinement depth under the same all-metric search
% as Figure~\ref{fig:supp-posthoc-token-dist}.}
% \label{fig:supp-posthoc-loop-dist}
% \end{figure}

% \bibliography{aaai2027}
% \end{document}


\documentclass[letterpaper]{article}
\usepackage[submission]{aaai2027}
\usepackage[hyphens]{url}
\usepackage{graphicx}
\urlstyle{rm}
\def\UrlFont{\rm}
\usepackage{natbib}
\usepackage{caption}
\frenchspacing
\usepackage{booktabs}
\usepackage{multirow}
\usepackage{amsmath}
\usepackage{amssymb}
\usepackage{xcolor}
\usepackage{array}

\newcommand{\ours}{ZipTok3D}
\pdfinfo{/TemplateVersion (2027.1)}

\renewcommand{\thesection}{S\arabic{section}}
\renewcommand{\thesubsection}{S\arabic{section}.\arabic{subsection}}
\renewcommand{\thefigure}{\Alph{figure}}
\renewcommand{\thetable}{\Alph{table}}
\renewcommand{\theequation}{S\arabic{equation}}
\setcounter{secnumdepth}{2}

\title{Supplementary Material for \ours{}: High-Fidelity 3D Tokenization with Compact Token Prefixes}
\author{Anonymous Submission}
\affiliations{}

\begin{document}
\maketitle

\section{Prefix-Length and Refinement-Depth Analysis}
\label{sec:full-sweeps}

\subsection{Complete Metric Sweeps}

% We evaluate prefix lengths
% $K\in\{1,2,4,8,16,32,64,128\}$ and refinement depths
% $L\in\{1,2,3,4,5,6\}$ using one selected checkpoint per dataset. At every
% operating point, the decoder receives a $K\times512$ latent tensor without
% suffix padding.

To characterize how prefix length and refinement depth jointly affect
reconstruction, we evaluate all supported prefix lengths
$K\in\{1,2,4,8,16,32,64,128\}$ and refinement depths
$L\in\{1,2,3,4,5,6\}$. This extends the query-IoU analysis in the main paper
to the complete operating range with two surface metrics. All operating
points within each dataset use the same selected checkpoint, and the decoder
receives the exact $K\times512$ prefix without suffix padding.

\begin{figure}[!t]
    \centering
    \captionsetup{skip=3pt}
    \begin{minipage}{\columnwidth}
        \centering
        \includegraphics[width=.90\columnwidth]
        {supplementary_figures/mesh_f1_interactions.pdf}
        \captionof{figure}{Mesh F1@0.02 across prefix lengths $K$ and
        refinement depths $L$ on ShapeNet (top) and TRELLIS (bottom).}
        \label{fig:supp-f1-curves}
    \end{minipage}
    \vspace{-0.4em}
    \begin{minipage}{\columnwidth}
        \centering
        \includegraphics[width=.90\columnwidth]
        {supplementary_figures/mesh_cd_interactions.pdf}
        \captionof{figure}{Mesh CD across prefix lengths $K$ and refinement
        depths $L$ on ShapeNet (top) and TRELLIS (bottom). Lower is better.}
        \label{fig:supp-cd-curves}
    \end{minipage}
\end{figure}

% Figures~\ref{fig:supp-f1-curves} and~\ref{fig:supp-cd-curves} extend the
% query-IoU analysis to the two surface metrics. Both exhibit the same asymmetric
% interaction: increasing $K$ is most beneficial under shallow decoding, whereas
% increasing $L$ is most beneficial for short prefixes. The
% differences between five and six passes are small at the reported operating
% points across query IoU, mesh CD, and mesh F1, supporting the use of five
% passes for evaluation.
Figures~\ref{fig:supp-f1-curves} and~\ref{fig:supp-cd-curves} provide
complementary axis-wise views of this interaction for mesh F1@0.02 and mesh
CD. When refinement is shallow, increasing $K$ yields the largest gains.
Conversely, when the prefix is short, increasing $L$ is most beneficial. This
asymmetric interaction is observed on both datasets and is consistent with the
query-IoU results in the main paper, showing that additional prefix tokens and
refinement passes contribute most strongly in different regimes.

Figure~\ref{fig:supp-heatmaps} reports the complete joint grids for query IoU,
mesh F1@0.02, and mesh CD. Across the evaluated range, the full sweeps confirm
that reconstruction is most sensitive to prefix length under shallow decoding,
whereas the benefit of additional refinement is concentrated at small $K$.
At the operating points used in the main paper, the differences between
$L=5$ and $L=6$ are small across all three metrics, supporting the use of five
refinement passes for evaluation.

\begin{figure*}[!t]
    \centering
    {\setlength{\tabcolsep}{1.5pt}
    \begin{tabular}{@{}cc@{}}
    \includegraphics[width=.45\textwidth]{supplementary_heatmaps_l1_l6/shapenet_query_iou.pdf} &
    \includegraphics[width=.45\textwidth]{supplementary_heatmaps_l1_l6/trellis_query_iou.pdf} \\[-.55em]
    \includegraphics[width=.45\textwidth]{supplementary_heatmaps_l1_l6/shapenet_mesh_f1_0p02.pdf} &
    \includegraphics[width=.45\textwidth]{supplementary_heatmaps_l1_l6/trellis_mesh_f1_0p02.pdf} \\[-.55em]
    \includegraphics[width=.45\textwidth]{supplementary_heatmaps_l1_l6/shapenet_mesh_cd.pdf} &
    \includegraphics[width=.45\textwidth]{supplementary_heatmaps_l1_l6/trellis_mesh_cd.pdf}
    \end{tabular}}
    \par\vspace{-.35em}
    \includegraphics[width=.85\textwidth]{supplementary_heatmaps_l1_l6/heatmap_relative_quality_legend.pdf}
    \par\vspace{-.25em}
    % \caption{Complete ShapeNet (left) and TRELLIS (right) sweeps for
    % refinement depths $L=1,\ldots,6$ and all supported prefix lengths. Rows
    % report query IoU, mesh F1@0.02, and mesh CD. Cell values are dataset means.
    % Colors are normalized independently within each dataset--metric panel. For
    % IoU and F1, error is $e=100-m$, where $m$ is the reported percentage; for
    % CD, $e=m$. Let $e_{\min}$ and $e_{\max}$ be the minimum and maximum errors
    % over all displayed $(K,L)$ cells in that panel. The normalized color score
    % is $q=\log(e_{\max}/e)/\log(e_{\max}/e_{\min})$, where $q=1$ denotes the
    % best cell and $q=0$ the worst. Darker colors indicate larger $q$.}
    \caption{Complete ShapeNet (left) and TRELLIS (right) sweeps for
    refinement depths $L=1,\ldots,6$ and all supported prefix lengths. Rows
    report query IoU, mesh F1@0.02, and mesh CD. Cell values are dataset means;
    CD is reported in units of $10^{-3}$. To orient all metrics so that better
    values receive darker colors, we define the residual from the optimum as
    $r=100-m$ for percentage-valued IoU and F1 and as $r=m$ for CD. Within each
    dataset--metric panel, we apply standard min--max normalization to
    $-\log r$. If $r_{\min}$ and $r_{\max}$ are the minimum and maximum
    residuals over its displayed $(K,L)$ cells, the resulting color score is
    $q=\log(r_{\max}/r)/\log(r_{\max}/r_{\min})$. Thus, $q=1$ denotes the best
    cell and $q=0$ the worst. Darker colors indicate larger $q$.}
    
    \label{fig:supp-heatmaps}
\end{figure*}

\subsection{Per-Object Refinement Behavior}

Dataset-level mean gains do not show whether refinement benefits most objects
or only a small subset. Table~\ref{tab:supp-refinement-diagnostics} therefore
reports the fraction of objects whose IoU, CD, or F1 improves over an early
refinement transition, $L=1\!\rightarrow\!3$, and a later transition,
$L=3\!\rightarrow\!5$, at two short-prefix settings,
$K\in\{1,4\}$.
% Mesh
% comparisons exclude seven
% TRELLIS objects for which at least one reconstructed mesh in the compared pair
% is empty or invalid, leaving 2,606 paired objects;
% all ShapeNet meshes are
% valid.
%注：写这些干什么，之后不要出现这个
%注：这段总体逻辑还可以。

\par\medskip
\noindent\begin{minipage}{\columnwidth}
\centering
\small
\setlength{\tabcolsep}{5pt}
\begin{tabular}{lccccc}
\toprule
& & & \multicolumn{3}{c}{Objects improved (\%)} \\
\cmidrule(lr){4-6}
Dataset & $K$ & Transition & IoU$\uparrow$ & CD$\downarrow$ & F1$\uparrow$ \\
\midrule
\multirow{4}{*}{ShapeNet}
 & 1 & $1\!\rightarrow\!3$ & 100.0 & 99.3 & 97.6 \\
 & 1 & $3\!\rightarrow\!5$ & 88.2 & 81.1 & 58.1 \\
 & 4 & $1\!\rightarrow\!3$ & 100.0 & 98.7 & 94.8 \\
 & 4 & $3\!\rightarrow\!5$ & 89.5 & 81.5 & 55.7 \\
\midrule
\multirow{4}{*}{TRELLIS}
 & 1 & $1\!\rightarrow\!3$ & 98.6 & 99.8 & 99.8 \\
 & 1 & $3\!\rightarrow\!5$ & 91.2 & 91.8 & 84.7 \\
 & 4 & $1\!\rightarrow\!3$ & 98.5 & 99.7 & 99.4 \\
 & 4 & $3\!\rightarrow\!5$ & 90.5 & 90.7 & 83.4 \\
\bottomrule
\end{tabular}
% \captionof{table}{Per-object improvement rates between selected refinement
% depths. Improvement is a strict increase for IoU and F1 and a strict decrease
% for CD; objects that tie or regress are not counted as improved. Values are
% rounded to one decimal place.}
\captionof{table}{Per-object improvement rates across selected refinement
intervals. An object is counted as improved if its IoU or F1 strictly increases
or its CD strictly decreases. Ties and regressions remain in the denominator
but are not counted as improvements. Values are rounded to one decimal place.}
\label{tab:supp-refinement-diagnostics}
\end{minipage}
\par\medskip

Across both datasets, increasing $L$ from 1 to 3 improves a large majority of
objects on every metric, indicating that the aggregate gains are not driven by
a small subset. Increasing $L$ further from 3 to 5 still improves more than
half of the objects in every setting, but the improvement rates are lower and
more metric-dependent, with the weakest consistency for ShapeNet mesh F1.
Early refinement therefore provides broadly shared benefits, whereas later
refinement remains useful for a smaller and more metric-dependent fraction of
objects.

\section{Experimental Setup Details}
\label{sec:experimental-details}

\subsection{Datasets, Preprocessing, and Splits}

\paragraph{ShapeNet.}
We use ShapeNetCore-v2 \citep{chang2015shapenet} with the 55-category
preprocessing and split manifests released by 3DShape2VecSet
\citep{zhang20233dshape2vecset} and adopted by COD-VAE
\citep{cho2025representing}. The Stage-1 tokenizer, prefix VAE, and
class-conditional EDM share the same partition: 48,597 objects are used for
training, 2,592 for validation and model selection, and a disjoint 1,283-object
test set for final evaluation. Reconstruction is evaluated on the complete test
set, whereas generation metrics use its category-specific subsets for airplane,
car, chair, table, and rifle. These role-based names correspond to the subsets
stored as \emph{test} and \emph{validation}, respectively, in the released
manifests.
For each object, the released data provide separate pools of uniform volume
queries and near-surface queries, their occupancy labels, a normalization
scale, and a dense reference surface point cloud. Following COD-VAE, we
consolidate the records by split while preserving category and object
identifiers. All manifest entries used in the reported experiments are
successfully converted. Query
coordinates are stored in FP16 for training and FP32 for both held-out splits,
while surface points remain in FP32.
During training, independent axis scales are sampled uniformly from
$[0.75,1.25]$, the shape is renormalized to the unit cube, and Gaussian noise
with standard deviation 0.005 is added to the surface points before clipping
to $[-1,1]$.  The same geometric scaling is applied to the occupancy queries.
Neither held-out split is augmented.

\paragraph{TRELLIS.}
TRELLIS-500K \citep{xiang2025structured} provides metadata for 500,777
training assets: 168,307 Objaverse-XL Sketchfab assets, 311,843 Objaverse-XL
GitHub assets, 4,485 ABO assets, 9,472 3D-FUTURE assets, and 6,670 HSSD assets.
Its separate 3,229-object Toys4K evaluation set is not used.  Our preprocessing
starts from the available downloaded assets rather than the complete metadata
catalogue.  We merge supported polygonal components and apply the
point-cloud-utils watertight conversion used by the VecSet preprocessing
pipeline at resolution 50,000 with seed 0.  The watertight mesh is centered at
its bounding-box center and scaled so that its maximum vertex radius is one.
Each successful object is represented by 100,000 surface points, 500,000
volume queries sampled uniformly from the radius-$\sqrt{3}$ sphere, and
500,000 near-surface queries divided equally between Gaussian perturbation
scales 0.005 and 0.05. Occupancy is one where the signed distance is
nonpositive and zero otherwise.
Per-object sampling is deterministic from seed 42 and the object identifier.
Inputs above 50,000 faces are decimated to that target before watertight
conversion when decimation succeeds. The preprocessing additionally sets a
900-second timeout and a 32-GiB memory limit per worker. An asset is excluded
if it lacks a supported,
nonempty polygonal mesh, contains nonfinite geometry, or fails mesh conversion,
normalization, surface sampling, signed-distance evaluation, or the stated
resource limits.

With seed 42, the successfully preprocessed assets are partitioned into 97\%
training, 2\% validation, and 1\% test subsets. The resulting test split
contains 2,613 objects: 57 ABO, 856 Objaverse-XL GitHub, and 1,700
Objaverse-XL Sketchfab assets. The anonymous code release provides the source
category and exact object identifier of every asset in this test split, which
is used for all reconstruction, bootstrap, and adaptive-budget analyses.

\subsection{Baseline Sources and Operating Points}

We compare with 3DILG \citep{zhang20223dilg}, VecSet
\citep{zhang20233dshape2vecset}, and COD-VAE
\citep{cho2025representing}.  These methods support deterministic
reconstruction at fixed, shape-independent token budgets.  All ShapeNet
comparisons use the common 1,283-object test split and the COD-VAE
metric implementation. The ShapeNet COD-VAE-32/64 checkpoints are inherited
from the official COD-VAE release, whereas COD-VAE-2 is trained using the
released implementation. On TRELLIS, COD-VAE-2/32/64 are initialized from
their corresponding ShapeNet checkpoints and fine-tuned on the common
preprocessed training pool. These checkpoints are selected by the highest mean
query IoU on the same validation split used by ZipTok3D.
VecSet-512 uses the inherited ShapeNet model without TRELLIS-specific training.
TRELLIS results are not reported for 3DILG or the shorter VecSet operating
points. All reported TRELLIS methods use the common test split and
evaluation protocol described above.

The reconstruction comparisons use 512 tokens for 3DILG; 32, 64, and 512
tokens for VecSet; and 2, 32, and 64 tokens for COD-VAE.  ZipTok3D is evaluated
with exact prefixes of 1, 2, and 4 tokens.  The class-conditioned generation
comparison uses the native stage-2 latent lengths of the baselines and the
exact two-token stage-2 representation of ZipTok3D.

\subsection{Stage-1 Optimization and Model Selection}
\label{sec:optimization}

For both datasets, each training example samples 2,048 points from the stored
surface pool as encoder input and separately samples 4,096 volume queries and
4,096 near-surface queries for reconstruction supervision. We optimize the
Stage-1 tokenizer with AdamW using $\beta_1=0.9$, $\beta_2=0.999$,
$\epsilon=10^{-8}$, weight decay 0.01, and an initial learning rate of
$10^{-4}$. Training uses FP16 mixed precision and clips the global gradient
norm at 0.5. Each of four A800 GPUs processes 56 examples, and gradients are
accumulated over three steps, giving an effective batch size of 672.

ShapeNet training runs for 1,000 epochs with a 50-epoch linear warmup followed
by cosine learning-rate decay. TRELLIS training runs for 300 epochs at a
constant learning rate. Every configuration is trained once with seed 123456.
For ShapeNet, checkpoints are selected by the highest mean query IoU on the
2,592-object validation split. For TRELLIS, the same criterion is evaluated on
the validation split. The implementation
uses Python 3.9, PyTorch 2.1.0, and CUDA 12.2. Training-time refinement and
supervision are specified in Section~\ref{sec:training-supervision}.

\subsection{Prefix VAE Optimization and Model Selection}

The prefix VAE uses AdamW with learning rate $10^{-4}$,
weight decay 0.01, global batch size 376, FP16 mixed precision, and gradient
clipping at 0.5.  It is trained for 100 epochs, with the learning rate halved at
epochs 60, 70, 80, and 90. It is trained on the same 48,597-object ShapeNet
training split as the Stage-1 tokenizer. Training uses four A800 GPUs and seed
123456. We select the checkpoint with the highest query IoU on the shared
2,592-object ShapeNet validation split.

\subsection{EDM Optimization and Model Selection}

The EDM uses AdamW with learning rate $10^{-4}$, weight decay 0.05, global
batch size 256, FP16 mixed precision, and gradient clipping at 1.0. Its
1,000-epoch schedule uses a 40-epoch linear warmup followed by cosine decay to
$10^{-6}$. Training uses four A800 GPUs and seed 123456. We set
$p_{\mathrm{mean}}=-1.2$, $p_{\mathrm{std}}=1.2$, and
$\sigma_{\mathrm{data}}=1$. The EDM is trained on latent arrays cached from the
shared ShapeNet training split, and its checkpoint is selected by the lowest
denoising loss on the shared validation split.

\section{Architecture and Objective Specifications}
\label{sec:architecture}

\subsection{Stage-1 Tokenizer Architecture}

Unless noted otherwise, the Stage-1 components follow the released COD-VAE design
\citep{cho2025representing}. ZipTok3D introduces nested prefix training, the
shared refinement block, and intermediate supervision.

The point encoder receives 2,048 surface points and forms 512 patches by
farthest-point sampling. It has four progressive blocks, each containing three
self-attention layers. The encoder width is 512, with eight attention heads,
a GEGLU MLP ratio of 4, and dropout 0.1. It produces a latent bank with maximum
length $M=128$ and width 512. The initial latent queries are the point
embeddings of 128 farthest-point-sampled surface positions, without a separate
learned latent-position table.

For each training example, one prefix length is sampled uniformly from
$\{1,2,4,8,16,32,64,128\}$. All suffix positions are masked from every decoder
attention operation. The decoder starts from 768 learned triplane tokens
$T^0$, corresponding to three 32-channel planes at spatial resolution
$128\times128$ with $8\times8$ patches. To predict occupancy, features are
bilinearly sampled from the three planes, summed, and processed by a
$32\!\rightarrow\!32\!\rightarrow\!1$ MLP with GELU.

\subsection{Triplane Initialization and Selection}

The selection block $\mathcal{S}_{\omega}$ first applies one cross-attention
layer between the 768 learned triplane tokens $T^0$ and the retained prefix
$Z_{:K}$. A LayerNorm--MLP uncertainty head assigns one score to each triplane
token. Hard descending top-$k$ selection retains the highest-scoring 25\%, or
192 tokens, as $H_K^0$. The remaining 576 tokens form the bypassed state
$R_K$ and retain their spatial indices for subsequent triplane restoration.
Selection is performed once before recurrent refinement rather than repeated
at every pass.

The ranking indices are discrete, so gradients do not propagate through the
top-$k$ ordering itself. Reconstruction gradients propagate through the
selected token features, while the uncertainty head receives direct
supervision from $\mathcal{L}_{\mathrm{unc}}$ defined below. The selected
indices are retained to restore the refined tokens to their original triplane
locations.

\subsection{Shared Refinement and Restoration}

The shared refinement module is a six-layer Transformer block of width 512
with eight attention heads, GEGLU MLP ratio 2, and dropout 0.1. It has no
causal mask, pass embedding, or pass-specific parameters. The recurrent
sequence at pass $\ell$ is
\begin{equation}
U_K^{\ell}=
\bigl[H_K^{\ell-1};Z_{:K}\bigr]
\in\mathbb{R}^{(192+K)\times512},
\end{equation}
where $H_K^{\ell-1}$ is the selected state and $Z_{:K}$ is the retained prefix.
The shared block processes this sequence, but only its first 192 outputs define
$H_K^{\ell}$. The prefix is reinserted unchanged at every pass, while the
bypassed tokens $R_K$ remain fixed outside the shared block.

After the final pass, the updated selected tokens are scattered to their
recorded spatial locations and combined with the 576 bypassed tokens at their
original locations. A linear residual head maps each restored spatial token to an
$8\times8\times32$ patch. The resulting residual is modulated by the predicted
uncertainty and added to the corresponding initial triplane patch.

\subsection{Auxiliary Losses and Intermediate Supervision}
\label{sec:training-supervision}

The two inherited COD-VAE auxiliary losses stabilize triplane initialization
and uncertainty-based selection. Let $a_j^0$ be the initial occupancy logit for
query $j$, $\hat y_j^0=\sigma(a_j^0)$ its probability, and
$\hat{\mathcal{Y}}^0=(\hat y_j^0)_{j=1}^{N_q}$ the corresponding probability
vector. The initialization loss applies the reconstruction criterion to this
initial prediction:
\begin{equation}
\mathcal{L}_{\mathrm{init}}
=\mathcal{L}_{\mathrm{rec}}(\hat{\mathcal{Y}}^0,\mathcal{Y}).
\end{equation}
The uncertainty target is the detached pointwise binary cross-entropy of the
initial prediction.  Following COD-VAE, it is clipped and normalized to
$[0,1]$ before regression:
\begin{align}
r_j & =
\frac{\operatorname{clip}(\operatorname{BCE}(\hat y_j^0,y_j)-0.01,0,0.99)}{0.99},\\
\mathcal{L}_{\mathrm{unc}}
&=\frac{1}{N_q}\sum_{j=1}^{N_q}(u_j-\operatorname{sg}(r_j))^2,
\end{align}
where $u_j$ is the uncertainty field decoded at the query. We use
$\lambda_{\mathrm{init}}=0.5$ and $\lambda_{\mathrm{unc}}=0.001$.

Stage-1 training unrolls six refinement passes. One intermediate depth is
sampled uniformly from $\{1,\ldots,5\}$ for each minibatch, while prefix
lengths are sampled independently for individual examples. The intermediate
term uses $\beta=0.5$, and its reconstruction and distillation components use
the same volume and near-surface weights, 1.0 and 0.1, as the final
reconstruction loss. Intermediate logits and detached final logits are
clamped to $[-30,30]$ before computing the distillation loss. The schedule
$\alpha_t$ is updated once per optimizer step and remains fixed across the
three microbatches accumulated into that update. Its endpoint $T_\alpha$ is
the scheduled optimizer-step count after distributed sharding, dropping
incomplete minibatches, and gradient accumulation. No additional
triplane-feature distillation term is used.

\subsection{Stage-2 Prefix VAE}

The prefix VAE accepts up to 16 ordered stage-1 tokens.  A LayerNorm and linear
projection independently map each 512-dimensional token to the mean and log
variance of a 32-dimensional diagonal Gaussian posterior.  The log variance is
clamped to $[-30,20]$, and training samples from the posterior by the
reparameterization trick.  A linear projection maps the sampled sequence back
to width 512, after which learned positional embeddings and a 12-layer
Transformer reconstruct the stage-1 token features.  This Transformer uses
eight heads, GEGLU MLP ratio 4, and dropout 0.1.  Its causal mask is
left-to-right:
\begin{equation}
A_{ij}=\begin{cases}
0, & j\leq i,\\
-\infty, & j>i,
\end{cases}
\end{equation}
so the reconstruction at position $i$ cannot use a later token.

Let $\widetilde Z$ denote a posterior sample and
$\bar Z=D_{\psi}(\widetilde Z)$ the reconstructed stage-1 token sequence.  The
stage-1 tokenizer is frozen while the prefix VAE is jointly optimized for
$\mathcal{B}=\{1,2,4,8,16\}$ with weights
$w_K\in\{1.2,1.1,1.0,0.9,0.8\}$.  Its objective is
\begin{equation}
\begin{aligned}
\mathcal{L}_{\mathrm{pVAE}}
&=\sum_{K\in\mathcal{B}}
\frac{w_K}{\sum_{K'\in\mathcal{B}}w_{K'}}
\\[-0.2em]
&\quad\times\Bigl[
\operatorname{MSE}\!\left(\bar Z_{:K},\operatorname{LN}(Z_{:K})\right)
\\[-0.2em]
&\hspace{3.2em}
+\mathcal{L}_{\mathrm{rec}}\!\left(F_6(\bar Z_{:K}),\mathcal{Y}\right)
\Bigr]
\\
&\quad+10^{-3}D_{\mathrm{KL}}\!\left(
q_{\psi}(\widetilde Z\mid Z_{:16})\,\|\,\mathcal{N}(0,I)
\right).
\end{aligned}
\end{equation}
where $F_6$ denotes the frozen stage-1 decoder and occupancy head evaluated
with six refinement passes.  The feature and occupancy terms both have unit
coefficients; the volume and near-surface weighting inside
$\mathcal{L}_{\mathrm{rec}}$ is given in
Section~\ref{sec:training-supervision}.
Causality ensures that every term for budget $K$ depends only on the
corresponding stage-2 prefix. The six-pass decoder is used only to supervise
prefix VAE training. During generation evaluation, sampled codes are decoded
using five refinement passes.

\subsection{Stage-2 Class-Conditional EDM}

We cache deterministic posterior means for the full $16\times32$ latent array,
normalize them using training-set channel statistics, and train the
class-conditional EDM on arrays from all 55 ShapeNet training categories.
Left-to-right causal self-attention ensures that the denoising prediction at
each position depends only on the corresponding noisy prefix. At the reported
two-token operating point, sampling therefore instantiates and denoises only
the first two latent positions. The EDM and prefix VAE decoder thus operate on
an exact $2\times32$ tensor without suffix padding.
The EDM follows the latent-array design used by
VecSet and COD-VAE \citep{zhang20233dshape2vecset,cho2025representing}: it has
width 384, 16 Transformer blocks, eight 48-dimensional heads, GEGLU MLP ratio
4, no dropout, learned token-position embeddings, and a left-to-right causal
self-attention mask. A
256-dimensional sinusoidal noise embedding is projected to width 384 and
provides scale and shift to the adaptive LayerNorms in every block.  The class
index spans all 55 ShapeNet categories. It is mapped to a learned
384-dimensional embedding and supplied as a one-token key and value context
through cross-attention in every block; token self-attention therefore models
the latent sequence while class conditioning is injected separately.

For a normalized latent array $x$, category $c$, and
$\log\sigma\sim\mathcal{N}(p_{\mathrm{mean}},p_{\mathrm{std}}^2)$, the EDM
objective is
\begin{equation}
\begin{aligned}
\mathcal{L}_{\mathrm{EDM}}
&=\mathbb{E}_{x,c,\sigma,\epsilon}\!\left[
\frac{\sigma^2+\sigma_{\mathrm{data}}^2}
{(\sigma\sigma_{\mathrm{data}})^2}
\right.\\[-0.2em]
&\hspace{4.2em}\left.
\times\left\|D_{\theta}(x+\epsilon;\sigma,c)-x\right\|_2^2
\right],
\\
&\hspace{4.2em}
\epsilon\sim\mathcal{N}(0,\sigma^2I).
\end{aligned}
\end{equation}
All latent positions present at an operating point are denoised in parallel;
there is no token-autoregressive sampling loop.

\section{Evaluation Protocols}
\label{sec:evaluation}

\subsection{Reconstruction Metric Settings}

For both datasets, the encoder receives 2,048 points sampled from the stored
surface pool of the target object. The task is therefore deterministic
reconstruction from sampled target surfaces rather than completion from
partial observations. For query IoU, the resulting representation is decoded
at all 500,000 stored volume-query locations per object.

For mesh evaluation, the decoder is queried on a dense $128^3$ occupancy grid.
We extract the zero-logit level set, equivalently the 0.5-probability
isosurface, with marching cubes. We then sample 100,000 points from the
reconstructed mesh and compare them with the complete stored reference surface
point cloud. Chamfer distance is the sum of the two mean bidirectional
Euclidean nearest-neighbor distances. At threshold 0.02, precision and recall
are the fractions of reconstructed and reference points, respectively, lying
within the threshold of the other surface. F1 is their harmonic mean.

\subsection{Class-Conditioned Generation Metric Settings}

Following COD-VAE \citep{cho2025representing}, we evaluate airplane, car,
chair, table, and rifle. We generate 2,000 shapes for each category and report
minimum matching distance (MMD-CD), coverage (COV-CD), and 1-nearest-neighbor
accuracy (1-NNA-CD). For a given category, the reference set $S_r$ contains all
objects of that category in the 1,283-object ShapeNet test split. From the 2,000
generated candidates, a deterministic subset of $5|S_r|$ shapes is used for
MMD-CD and COV-CD, while a subset of $|S_r|$ shapes is used for 1-NNA-CD. The
same reference sets and generated-set construction are used for every method.

Each mesh is represented by 2,048 points sampled from its surface. Using CD as
the set distance, MMD averages the distance from each reference shape to its
nearest generated shape. COV is the fraction of reference shapes selected as a
nearest neighbor of at least one generated shape. The fixed 5:1
generated-to-reference ratio makes MMD-CD and COV-CD comparable across methods.
1-NNA is the leave-one-out nearest-neighbor classification accuracy on the
union of equally sized reference and generated sets, avoiding sample-count
imbalance between the two sets. We report the unweighted mean of each metric
over the five categories.

VecSet and COD-VAE use 18 diffusion steps, whereas 3DILG uses its native
512-step autoregressive sampler.  Our EDM uses 18 Heun steps with $\rho=7$,
$\sigma_{\max}=80$, and $\sigma_{\min}=0.002$, without condition dropout or
classifier-free guidance.

\subsection{Efficiency Measurement Protocol}

All efficiency measurements use the corresponding trained checkpoints on one
NVIDIA H20 GPU with batch size 16 and FP32 inference.  TF32 and Flash SDP are
disabled.  CUDA events measure GPU execution after warmup; model loading,
data loading, CPU--GPU transfers, marching cubes, and file output are excluded.
For reconstruction, full inference begins with the input point cloud and
includes point encoding, tokenizer decoding, and all $128^3$ occupancy queries
in chunks of 131,072 points. Decoder latency instead begins with an already
available stage-1 prefix and ends when the complete triplane has been
produced; it excludes point encoding and occupancy queries.

For generation, sampling throughput measures EDM latent generation alone.
Full throughput begins from the same latent noise and additionally includes
prefix VAE decoding, five-pass stage-1 tokenizer decoding, and the complete
$128^3$ occupancy-field query. Thus the point encoder is part of reconstruction
full inference but not of the generation pipeline. We discard three warmup
batches and average ten measured batches for reconstruction and decoder
latency; generation uses five warmup and thirty measured batches. For a
measured mean batch latency $t$ in milliseconds,
\begin{equation}
\mathrm{throughput}=\frac{16\times1000}{t}.
\end{equation}
Peak allocated CUDA memory is measured after resetting the peak-memory counter
following warmup. It is measured over the corresponding full pipeline;
reserved memory is not used.

\subsection{Reconstruction Efficiency Results}

\begin{table}[t]
\centering
\small
\setlength{\tabcolsep}{4pt}
\begin{tabular}{lrrr}
\toprule
Method & Tok. & Full shape/s & Peak GiB \\
\midrule
VecSet & 512 & 4.20 & 32.63 \\
COD-VAE & 2 & 82.48 & 2.08 \\
COD-VAE & 32 & 79.82 & 2.08 \\
COD-VAE & 64 & 77.29 & 2.09 \\
\midrule
\ours{} & 1 & 72.51 & 2.03 \\
\ours{} & 2 & 72.51 & 2.03 \\
\ours{} & 4 & 72.24 & 2.03 \\
\bottomrule
\end{tabular}
\caption{Full reconstruction efficiency under the common protocol.  ZipTok3D
uses five refinement passes.}
\label{tab:supp-reconstruction-efficiency}
\end{table}

Table~\ref{tab:supp-reconstruction-efficiency} reports the resulting
representation--compute tradeoff. Shorter prefixes reduce neither the dense
occupancy query nor the number of shared refinement calls. Consequently,
ZipTok3D has lower end-to-end throughput than the single-pass COD-VAE variants,
despite its shorter latent sequence. Peak memory remains similar because the
common encoder and dense-field query dominate the allocation.

\subsection{Statistical Analysis}

For the primary comparison between ZipTok3D at $K=4,L=5$ and COD-VAE-32, we
resample aligned TRELLIS objects with replacement. Each of 20,000 resamples
contains 2,613 object indices. For each metric, we compute the difference
between the paired sample means and report the 2.5th and 97.5th percentiles of
the resulting bootstrap distribution. The bootstrap uses seed 123456, and all
effects are defined as ZipTok3D minus COD-VAE, so negative CD is favorable.
These intervals quantify variation over evaluation objects conditional on the
trained models; they do not measure variation across independent training
runs.


\section{Post-Hoc Adaptive-Budget Diagnostic}
\label{sec:posthoc-oracle}

\begin{table}[t]
\centering
{\small\setlength{\tabcolsep}{2pt}
\begin{tabular*}{\columnwidth}{@{\extracolsep{\fill}}llcccc@{}}
\toprule
Data & Method & Found & Rate & Avg. $K$ & Avg. $L$ \\
\midrule
ShapeNet & COD-VAE-32 & -- & -- & 32.00 & -- \\
ShapeNet & \ours{} & 979/1,283 & 76.31 & 2.43 & 3.41 \\
TRELLIS & COD-VAE-32 & -- & -- & 32.00 & -- \\
TRELLIS & \ours{} & 1,767/2,613 & 67.62 & 3.02 & 3.97 \\
\bottomrule
\end{tabular*}
\vspace{0.35em}

\begin{tabular*}{\columnwidth}{@{\extracolsep{\fill}}llccc@{}}
\toprule
Data & Method & IoU$\uparrow$ & CD$\downarrow$ & F1$\uparrow$ \\
\midrule
ShapeNet & COD-VAE-32 & 97.2 & 0.011 & 98.1 \\
ShapeNet & \ours{} & 97.3 & 0.011 & 98.3 \\
TRELLIS & COD-VAE-32 & 75.8 & 0.015 & 96.1 \\
TRELLIS & \ours{} & 76.6 & 0.014 & 96.6 \\
\bottomrule
\end{tabular*}}
\caption{All-metric post-hoc diagnostic. The \ours{} rows report the post-hoc
oracle. Found and Rate give the number and percentage of shapes for which the
search identifies an operating point that matches or exceeds COD-VAE-32 on all
three rounded per-shape metrics. Avg. $K$ and Avg. $L$ are computed over these
shapes. Metrics in the lower block are aggregated over the same subset for
both methods, so the COD-VAE rows are subset rather than full-split results.
IoU and F1 are percentages.}
\label{tab:supp-posthoc-all}
\end{table}



\begin{figure}[!t]
    \centering
    \captionsetup{skip=3pt}
    \begin{minipage}{\columnwidth}
        \centering
        \includegraphics[width=.94\columnwidth]
        {supplementary_figures/adaptive_success_token_distribution.pdf}
        \captionof{figure}{Distribution of the first successful prefix length
        $K$ under the token-first all-metric search. Percentages are normalized
        within the successful subset of each dataset.}
        \label{fig:supp-posthoc-token-dist}
    \end{minipage}
    \vspace{-0.35em}
    \begin{minipage}{\columnwidth}
        \centering
        \includegraphics[width=.94\columnwidth]
        {supplementary_figures/adaptive_success_loop_distribution.pdf}
        \captionof{figure}{Distribution of the first successful refinement
        depth $L$ under the same search and normalization.}
        \label{fig:supp-posthoc-loop-dist}
    \end{minipage}
\end{figure}

To characterize how the required budget varies across individual shapes, we
perform a post-hoc diagnostic over
$K\in\{1,2,4,8,16,32\}$ and $L\in\{1,\ldots,6\}$ from the same checkpoint.
Candidates are searched by increasing $K$ and then increasing $L$, prioritizing
a shorter representation before shallower decoding.  This diagnostic accesses
ground-truth reconstruction metrics and thus measures adaptive potential rather
than the performance of a deployable budget predictor.

We use a strict criterion that accepts a candidate only when
all three metrics are no worse than COD-VAE-32 after decimal half-up rounding
(one decimal for IoU/F1 and three decimals for CD).  Shapes without a successful
candidate are excluded from this diagnostic.  Table~\ref{tab:supp-posthoc-all}
reports both methods on the same successful subset, preventing differences in
subset composition from affecting the comparison.


No candidate in the evaluated grid satisfies the strict all-metric criterion
for the remaining 304 ShapeNet and 846 TRELLIS objects. All fixed-budget
evaluations include these objects; they are excluded only from the oracle
averages in Table~\ref{tab:supp-posthoc-all}. Consequently, the
reported average budgets characterize the successful subset rather than the
expected cost of a deployable early-exit policy.

Figures~\ref{fig:supp-posthoc-token-dist}
and~\ref{fig:supp-posthoc-loop-dist} summarize the first successful budgets
within this subset. Three or four passes account for 68.3\% and 71.6\% of the
ShapeNet and TRELLIS subsets, respectively. This distribution is consistent
with the dataset-level saturation behavior in
Figures~\ref{fig:supp-f1-curves} and~\ref{fig:supp-cd-curves}.

\makeatletter
\setlength{\@fptop}{0pt}
\makeatother

\begin{figure}[!t]
\centering
\includegraphics[width=.95\columnwidth]
{supplementary_figures/additional_qualitative_comparisons.png}
\captionsetup{width=\columnwidth,skip=3pt}
\caption{Additional reconstruction comparisons on ShapeNet (top five rows)
and TRELLIS (bottom two rows). ZipTok3D uses $K=4$ and $L=5$; the baseline
uses 32 COD-VAE tokens.}
\label{fig:supp-additional-qualitative}
\end{figure}

\section{Additional Qualitative Comparisons}

Figure~\ref{fig:supp-additional-qualitative} extends the qualitative comparison
to objects with thin components, large openings, curved surfaces, and repeated
architectural elements. The open shelf, chair, and table retain their principal
empty regions, while the vessel and aircraft preserve separated components and
overall topology. On TRELLIS, both reconstructions recover the multi-level
building layouts and repeated roof or fortification structures. The remaining
differences are concentrated around small protrusions, thin supports, roof
eaves, and crenellations, which remain challenging under both token budgets.



\clearpage

\makeatletter
\setlength{\@fptop}{0pt plus 1fil}
\makeatother

\bibliography{aaai2027}
\end{document}
